\documentclass[10pt,a4paper]{article}
\usepackage{fontspec,amsmath,amssymb,graphicx}
\usepackage[margin=18mm]{geometry}
\setmainfont{Times New Roman}
\setlength{\parindent}{0pt}
\setlength{\parskip}{5pt}
\allowdisplaybreaks
\sloppy
\emergencystretch=2em
\begin{document}
\section*{VD\_4\_6\_2\_03 — Lời giải bổ sung}
\subsection*{Kiến thức và trình tự giải}
\textbf{Bước 1 — Chọn nhân nhiệt theo miền và hệ số.} Trên $\mathbb R$, $g_K(z,t)=(4\pi Kt)^{-1/2}e^{-z^2/(4Kt)}$, $u=g_K*f$. Trên nửa trục, thác triển lẻ cho Dirichlet $0$, chẵn cho Neumann $0$; sau đổi $y\mapsto-y$ ở nửa âm:


\[
H_D(x,y,t)=g_K(x-y,t)-g_K(x+y,t),\qquad
H_N(x,y,t)=g_K(x-y,t)+g_K(x+y,t).
\]


Kiểm tra $H_D(0,y,t)=0$, $\partial_xH_N(0,y,t)=0$. Nửa trục phải có một điều kiện biên; dữ liệu đầu một mình không xác định duy nhất.

\textbf{Bước 2 — Tính các tích phân không gian.} Mode $\sin ky,\cos ky$ biến thành $e^{-Kk^2t}\sin kx,e^{-Kk^2t}\cos kx$. Với Gauss $e^{-\alpha(y-m)^2}$, $\alpha>0$, hoàn thành bình phương:


\[
-\frac{(x-y)^2}{4Kt}-\alpha(y-m)^2
=-\frac{1+4\alpha Kt}{4Kt}\left(y-\frac{x+4\alpha Ktm}{1+4\alpha Kt}\right)^2
-\frac{\alpha(x-m)^2}{1+4\alpha Kt}.
\]


Đổi biến để $\int_{\mathbb R}e^{-z^2}dz=\sqrt\pi$, suy $(g_K*f)(x)=e^{-\alpha(x-m)^2/(1+4\alpha Kt)}/\sqrt{1+4\alpha Kt}$. Với dữ liệu chỉ thị đoạn, cận hữu hạn tạo $\operatorname{erf}$; với $e^{-y}$ trên nửa trục, cận $0,\infty$ tạo $\operatorname{erfc}$, trong đó $\operatorname{erf}z=2\int_0^ze^{-r^2}dr/\sqrt\pi$, $\operatorname{erfc}z=1-\operatorname{erf}z$.

\textbf{Bước 3 — Thêm nguồn và biên.} Nguồn cho $u=g_K*f+\int_0^tg_K(\cdot,t-s)*F(\cdot,s)\,ds$, hoặc thay $g_K$ bằng $H_D,H_N$. Biên Dirichlet $\mu(t)$ trên nửa trục thêm


\[
\int_0^t\frac{x\,e^{-x^2/[4K(t-s)]}}{\sqrt{4\pi K}(t-s)^{3/2}}\mu(s)ds.
\]


Nguồn một mode giảm cùng tốc độ khuếch tán cho tích phân thời gian hằng sau khi gom mũ, nên xuất hiện $te^{-Kk^2t}$.

\textbf{Bước 4 — Đổi biến khi hệ số phụ thuộc thời gian hoặc có vận chuyển.} Với $u_t=a(t)u_{xx}+b(t)u_x+c(t)u$, đặt $A=\int_0^ta$, $B=\int_0^tb$, $C=\int_0^tc$, $u=e^Cv(x+B,A)$. Tính $u_t=e^C(cv+bv_z+av_\tau)$, $u_x=e^Cv_z$, $u_{xx}=e^Cv_{zz}$, suy $v_\tau=v_{zz}$. Dấu dịch là $x+B$, không phải $x-B$. Với hạng $xu_x$, dùng $z=e^{\pm t}x$ và chọn đồng hồ $\tau'(t)=e^{\pm2t}$.

\textbf{Bước 5 — Nhiều chiều và kiểm tra.} Nhân $n$ chiều là tích các nhân một chiều. Dữ liệu dạng tích tách được tích phân không gian bằng Fubini; nguồn dạng tích vẫn phải tích phân theo một thời gian chung $s$. Kiểm tra PDE bằng đạo hàm Duhamel: đạo hàm cận trên trả lại $F(x,t)$, phần còn lại bằng $K\Delta u$. Nhân Gauss là xấp xỉ đơn vị, nên trả lại dữ liệu tại điểm liên tục; dữ liệu nhảy không cho nghiệm liên tục đến mọi điểm lúc đầu.
\subsection*{Lời giải áp dụng cho bài}
\textbf{Bước 1. Chuyển bài $u_t=u_{xx}+xu$ về nhiệt.}
Đặt $u=e^{tx+t^3/3}v(z,t)$, $z=x+t^2$. Tính
\begin{align*}
u_t&=e^{tx+t^3/3}\{(x+t^2)v+2tv_z+v_t\},\\
u_{xx}+xu&=e^{tx+t^3/3}\{(t^2+x)v+2tv_z+v_{zz}\}.
\end{align*}
Do đó $v_t=v_{zz}$, $v(z,0)=f(z)$. Phép đổi biến này biện minh công thức cho $f$ liên tục bị chặn mà không cần đánh giá $\widehat f$ tại tần số phức.
\textbf{Bước 2. Fourier cho bài đã đổi.} Với quy ước $\widehat v=\int e^{-i\xi z}v\,dz$, $\widehat v_t=-\xi^2\widehat v$, suy $\widehat v=e^{-t\xi^2}\widehat f$. Biến đổi ngược cho $v=g_1*f$. Vì vậy
\[u=e^{tx+t^3/3}\int_{\mathbb R}g_1(x+t^2-y,t)f(y)dy.\]
Gom mũ:
\[tx+t^3/3-\frac{(x+t^2-y)^2}{4t}
=-\frac{(x-y)^2}{4t}+\frac{t(x+y)}2+\frac{t^3}{12}.\]
Vậy
\[u=\int_{\mathbb R}G(x,y,t)f(y)dy,\qquad
G=\frac1{\sqrt{4\pi t}}
\exp\left[-\frac{(x-y)^2}{4t}+\frac{t(x+y)}2+\frac{t^3}{12}\right].\]
\textbf{Bước 3. Kiểm tra Green.} $G>0$, đối xứng theo $x,y$, trơn với $t>0$. Với $\ell=\log G$,
\[\ell_x=-\frac{x-y}{2t}+\frac t2,\quad \ell_{xx}=-\frac1{2t},\quad
\ell_t=-\frac1{2t}+\frac{(x-y)^2}{4t^2}+\frac{x+y}2+\frac{t^2}4.\]
Ta có $\ell_t=\ell_{xx}+\ell_x^2+x$, suy $G_t=G_{xx}+xG$. Từ biểu diễn Gauss đã dịch,
\[\int_{\mathbb R}G(x,y,t)dy=e^{tx+t^3/3}.\]
Nhân không có tổng bằng $1$ do hạng phản ứng $xu$.
\textbf{Bước 4. Dữ liệu đầu và sự tập trung.} Với $x$ trong tập compact, hệ số $e^{tx+t^3/3}\to1$ và tâm $x+t^2\to x$. Gauss là xấp xỉ đơn vị, nên $u(x,t)\to f(x)$ tại điểm liên tục, cũng liên tục khi $(x,t)\to(x_0,0)$. Với $\delta>0$, tích phân nhân trên $|y-x|>\delta$ tiến về $0$. Điều kiện này phải là khoảng cách $|y-x|$, không phải $|x|$ như một dòng sai biến trong bản in. Công thức không bảo đảm nghiệm bị chặn theo $x$ khi $t>0$, vì phản ứng tuyến tính theo $x$ có thể làm nghiệm tăng theo mũ.\end{document}